\section{Guarantees and Correctness Properties}\label{sec:properties}

We now formally state the three correctness properties guaranteed by the Bailout Protocol: \SAFETY{} (no autonomous collapse), \LIVENESS{} (eventual human contact), and \ACCOUNTABILITY{} (human always identifiable). For each property we give a formal statement, the structural assumptions under which it holds, and a proof sketch that traces the argument through the protocol's specification. The proofs are constructive: they show how each property follows directly from the deterministic state machine definition, the bypass mechanism, and the timeout architecture.

% --------------------------------------------------------------
\subsection{Overview of Correctness Properties}
% --------------------------------------------------------------

The Bailout Protocol is designed to satisfy three mutually independent correctness properties that together establish the upstream accountability guarantee. These properties are not aspirational — they are derived from the protocol's formal specification and hold under the stated structural assumptions.

\begin{enumerate}
\item[\textbf{SAFETY.}] No \bxthree{} node may remain in a state of autonomous operation after entering the Bailout Protocol. Entry to \texttt{ESCALATING} is atomic and irreversible by any machine actor; the protocol terminates only at $h(n)$.

\item[\textbf{LIVENESS.}]. Every exception that satisfies the trigger condition (TC) eventually reaches the human anchor $h(n)$, provided the deployment satisfies its provisioned timeout bounds. There is no execution path in which an unresolvable exception is permanently suppressed.

\item[\textbf{ACCOUNTABILITY.}] Every directive recorded in the Forensic Ledger at \texttt{RESOLVED} is attributable to a named human principal. No machine actor can issue, simulate, or substitute a directive at the \texttt{RESOLVED} state.
\end{enumerate}

These properties are inter-referential but not redundant. \SAFETY{} constrains what may happen during the protocol run; \LIVENESS{} constrains what must eventually happen; \ACCOUNTABILITY{} constrains who the terminal actor must be. Together they eliminate the three structural failure modes identified in \S~\ref{sec:bailout}: autonomous absorption, silent suppression, and anonymous delegation.

% --------------------------------------------------------------
\subsection{SAFETY: No Autonomous Collapse}
% --------------------------------------------------------------

\subsubsection{Formal Statement.}

\begin{bluebox}
\textbf{Theorem (Safety).} Under Assumptions (A1)--(A5), for every node $n \in \text{nodes}$ and every reachable state $q$ of $\mathcal{B}$:

\[
\forall n \in \text{nodes},\; \forall q \in Q,\; q = \text{ESCALATING} \;\implies\; \text{terminus}(q) = \{h(n)\}.
\tag{SAFETY}
\]

Equivalently: no \bxthree{} node in \texttt{ESCALATING} state may transition to a state whose termination set contains any machine actor. The protocol has no machine-actor terminal state.
\end{bluebox}

\subsubsection{Structural Assumptions.}

\begin{enumerate}
\item[(A1)] The parent pointer $\alpha(n)$ and human-root identifier $h(n)$ are immutable at the machine-actor level; they are stored in \fact-layer-managed memory and cannot be modified by any machine actor at runtime.
\item[(A2)] The \fact{} layer pre-commit hook rejects any ledger entry that would retroactively authorize a resolution after $T_{\text{bounds}}$ has expired without a prior \fact-layer-approved resolution.
\item[(A3)] The state transition relation $\rightarrow$ is deterministic: for every $(q, \sigma) \in Q \times \Sigma$, at most one $q' \in Q$ satisfies $(q, \sigma, q') \in \rightarrow$.
\item[(A4)] Entry to \texttt{ESCALATING} is atomic: the \fact{} layer executes \texttt{SendToParent} immediately upon the T4 transition, with no intervening machine-actor decision step.
\item[(A5)] The set of terminal accepting states $Q_F = \{\text{RESOLVED}\}$ contains no state that permits autonomous continuation.
\end{enumerate}

\subsubsection{Proof Sketch.}

The argument proceeds by exhaustively examining the state machine's transition relation and showing that \texttt{ESCALATING} has no outgoing transition that leads to any state controlled by a machine actor.

\medskip
\noindent\textbf{Step 1: Characterize the terminal states of $\mathcal{B}$.} By definition of $Q_F$, the only terminal accepting state is \texttt{RESOLVED}. The state definitions (\S~\ref{sec:state-machine}) establish that \texttt{RESOLVED} is reached only upon receipt of a directive $d \in \{\texttt{ACK}, \texttt{RESOLVE}, \texttt{REJECT}\}$ issued by $h(n)$. No machine actor may issue any of these directives by architectural constraint of the bypass mechanism (BP, \S~\ref{sec:bypass}). Hence any path to \texttt{RESOLVED} necessarily passes through $h(n)$.

\medskip
\noindent\textbf{Step 2: Show that \texttt{ESCALATING} has no machine-actor outgoing transition.} The transition table (Table~\ref{tab:bailout-transition}) defines the outgoing transitions from \texttt{ESCALATING} as follows:
\begin{itemize}
\item $e_{\text{auth}}$ (T6): transitions to \texttt{HUMAN\_ACKNOWLEDGED}. This event is issued by the \fact{} layer upon verification that the acknowledgment originated from $h(n)$, not from any machine actor.
\item $e_{\text{ack}}$ (T5): transitions to \texttt{IDLE}. This is the case where the human anchor $h(n)$ reviews the exception and determines no further action is required.
\item $e_{\text{resolve}}$ or $e_{\text{reject}}$ (T7): transitions to \texttt{RESOLVED}. Both require a binding directive from $h(n)$.
\end{itemize}
No transition from \texttt{ESCALATING} is triggered by any machine-actor event. The only possible destinations — \texttt{HUMAN\_ACKNOWLEDGED}, \texttt{IDLE}, and \texttt{RESOLVED} — all require human-actor involvement.

\medskip
\noindent\textbf{Step 3: Show atomicity of entry to \texttt{ESCALATING}.} By (A4) and the T4 transition definition (\S~\ref{sec:transitions}), entry to \texttt{ESCALATING} occurs via an immediate atomic transition from \texttt{BAIL\_PENDING} upon receipt of the compounded timeout signal $(T_{\text{bounds}}, T_{\text{prop}})$. The \fact{} layer executes \texttt{SendToParent} directly; no machine-actor decision gate exists between \texttt{BAIL\_PENDING} and \texttt{ESCALATING}. By INV-2 (\S~\ref{sec:state-machine}), this transition is the \emph{only} permissible transition from \texttt{BAIL\_PENDING}. No machine actor can suppress, delay, or redirect this transition.

\medskip
\noindent\textbf{Step 4: Combine via the bypass property.} By (A1) and the bypass mechanism (\S~\ref{sec:bypass}), the communication path from \texttt{ESCALATING} to $h(n)$ contains no machine actor in the termination set. The \texttt{SendToParent} call routes directly to the human anchor via the provisioned pathway; intermediate machine nodes forward without attempting local resolution. By BP, $\text{terminus}(\text{ESCALATING}(n)) = \{h(n)\}$.

\medskip
\noindent\textbf{Conclusion.} Combining Steps 1--4, any node that enters \texttt{ESCALATING} must eventually reach \texttt{RESOLVED}, and any path to \texttt{RESOLVED} requires a directive from $h(n)$. No machine actor can absorb, redirect, or suppress the escalation. $\square$

% --------------------------------------------------------------
\subsection{LIVENESS: Eventual Human Contact}
% --------------------------------------------------------------

\subsubsection{Formal Statement.}

\begin{bluebox}
\textbf{Theorem (Liveness).} Under Assumptions (L1)--(L3), for every node $n \in \text{nodes}$ and every exception $x$ satisfying the trigger condition $\text{TRIGGER}(n, x)$, the protocol run from the firing of T1 reaches $h(n)$ within bounded wall-clock time $T_{\max}$:

\[
\forall n \in \text{nodes},\; \forall x,\; \text{TRIGGER}(n, x) \;\implies\; \exists t \leq T_{\max}:\; \text{received}(h(n), e, t).
\tag{LIVENESS}
\]

where $T_{\max} = T_{\text{bounds}} + N_{\max} \cdot T_{\text{cap}} + T_{\text{human}}$.
\end{bluebox}

\subsubsection{Structural Assumptions.}

\begin{enumerate}
\item[(L1)] The Pathway Health Registry reports at least one functional communication pathway to $h(n)$ within any sliding window of duration $T_{\text{prop}}$.
\item[(L2)] $h(n)$ is provisioned with at least one communication pathway and responds to \texttt{ACK} or directive messages within $T_{\text{human}}$ with probability 1.
\item[(L3)] The escalation algorithm retries at most $N_{\max}$ times, each retry bounded by $T_{\text{cap}}$; after $N_{\max}$ retries, the protocol reaches \texttt{RESOLVED} with \texttt{PATHWAY\_EXHAUSTED}.
\end{enumerate}

Assumption (L1) is a deployment requirement: it is the responsibility of the deploying organization to provision sufficient communication redundancy to satisfy it. Assumptions (L2) and (L3) are protocol invariants enforced by the \fact{} layer.

\subsubsection{Proof Sketch.}

The liveness proof proceeds by bounded case analysis over the possible execution paths of the escalation algorithm and the state machine.

\medskip
\noindent\textbf{Step 1: Trigger to \texttt{ESCALATING} is bounded.} By T1, the trigger condition causes an immediate transition to \texttt{BOUNDED}. By T3, if no \fact-layer-approved resolution arrives within $T_{\text{bounds}}$, the machine transitions to \texttt{BAIL\_PENDING}. By T4, entry to \texttt{ESCALATING} is atomic and immediate upon the compounded timeout signal. The elapsed time from T1 to \texttt{ESCALATING} entry is therefore exactly $T_{\text{bounds}}$.

\medskip
\noindent\textbf{Step 2: Each retry iteration is bounded by $T_{\text{cap}}$.} The escalation algorithm executes the loop:

\begin{trivlist}\itemsep=0pt
\item[]\textbf{while} $retry\_count < N_{\max}$ \textbf{do}
\item[]\hspace{1em}$ack\_received \leftarrow \texttt{SendToParent}(n, e, C, pathway, T_{\text{current}})$
\item[]\hspace{1em}\textbf{if} $ack\_received$ \textbf{then return}
\item[]\hspace{1em}$retry\_count \leftarrow retry\_count + 1$
\item[]\hspace{1em}$T_{\text{current}} \leftarrow \min(2 \cdot T_{\text{current}},\; T_{\text{cap}})$
\end{trivlist}

Each invocation of \texttt{SendToParent} waits at most its allocated timeout $T_{\text{current}}$, which starts at $T_{\text{prop}}$ and doubles on each retry until it reaches $T_{\text{cap}}$. Since $T_{\text{current}}$ is monotonically non-decreasing and bounded above by $T_{\text{cap}}$, each iteration completes within at most $T_{\text{cap}}$ wall-clock time. The loop executes at most $N_{\max}$ times. The total elapsed time for the loop is therefore at most $N_{\max} \cdot T_{\text{cap}}$.

\medskip
\noindent\textbf{Step 3: Human acknowledgment is bounded by $T_{\text{human}}$.} Upon the first successful \texttt{SendToParent} delivery to $h(n)$, the state machine transitions to \texttt{HUMAN\_ACKNOWLEDGED} via T6. By (L2), $h(n)$ issues a directive within $T_{\text{human}}$. This produces either T5 (return to \texttt{IDLE}) or T9 (enter \texttt{RESOLVED}), both of which are terminal for the protocol run.

\medskip
\noindent\textbf{Step 4: Pathway exhaustion is terminal but still recorded.} If all $N_{\max}$ retries fail to deliver to $h(n)$ (the Pathway Health Registry reports no functional pathway, or all pathways fail $K$ consecutive health checks), the algorithm exits the loop and fires \texttt{PATHWAY\_EXHAUSTED}. The protocol transitions to \texttt{RESOLVED} with this tag, satisfying the liveness bound $T_{\max}$ from below: the protocol has reached a terminal state with a Forensic Ledger record, and the \texttt{PATHWAY\_EXHAUSTED} tag signals the deploying organization that human re-engagement is required to restore pathway connectivity.

\medskip
\noindent\textbf{Step 5: Combine the bounds.} The total elapsed wall-clock time from T1 firing to terminal state is:

\[
T_{\text{bounds}} + N_{\max} \cdot T_{\text{cap}} + T_{\text{human}} = T_{\max}.
\]

This bound holds for all possible execution paths. No execution path can exceed $T_{\max}$ because: (a) the trigger-to-escalation path is fixed at $T_{\text{bounds}}$; (b) the retry loop is bounded by $N_{\max} \cdot T_{\text{cap}}$ by construction; (c) the human response loop is bounded by $T_{\text{human}}$ by (L2). $\square$

\medskip
\noindent\textbf{Remark.} The liveness bound $T_{\max}$ is a deployment parameter, not a protocol constant. A deployment in a high-latency environment (e.g., satellite-linked agricultural sensors in the San Luis Valley) must provision timeouts accordingly and certify that $T_{\max}$ remains within the safety threshold for the node's operational domain.

% --------------------------------------------------------------
\subsection{ACCOUNTABILITY: Human Always Identifiable}
% --------------------------------------------------------------

\subsubsection{Formal Statement.}

\begin{bluebox}
\textbf{Theorem (Accountability).} Under Assumptions (C1)--(C4), for every terminal protocol run reaching \texttt{RESOLVED}, the Forensic Ledger entry for that run contains a non-anonymous reference to a named human principal $h(n)$ who issued the binding directive:

\[
\forall n \in \text{nodes},\; \forall \text{run}_\rho \text{ of } \mathcal{B},\; 
\text{state}(\text{run}_\rho) = \text{RESOLVED} \;\implies\; 
\exists!\, h(n) \in \text{HumanPrincipals}:\; \text{directive}(\text{run}_\rho) \in \text{Directives}(h(n)).
\tag{ACCOUNTABILITY}
\]

Furthermore, no machine actor may appear in the attribution chain between the originating node and $h(n)$.
\end{bluebox}

\subsubsection{Structural Assumptions.}

\begin{enumerate}
\item[(C1)] Every human principal $h(n)$ is identified by a unique, non-repudiable credential established at node provisioning. This credential is stored in the Fact Layer authorization ledger and cannot be altered at runtime.
\item[(C2)] The \fact{} layer verifies the provenance of every directive at \texttt{RESOLVED} by cryptographic attestation of the human principal's credential. A directive that cannot be attributed to a verified human principal is rejected and logged as a protocol violation.
\item[(C3)] The \purpose{} layer is non-delegable as an exception terminus: no machine actor may receive an exception and issue a binding directive that substitutes for a human directive.
\item[(C4)] The Forensic Ledger is append-only and tamper-evident: any attempt to modify or delete a \texttt{RESOLVED} entry after the fact is detectable by the integrity hash chain.
\end{enumerate}

\subsubsection{Proof Sketch.}

The accountability proof has two components: (a) identification — $h(n)$ is the only entity capable of issuing the terminal directive; (b) attribution — the directive is recorded with $h(n)$'s non-repudiable credential in the Forensic Ledger.

\medskip
\noindent\textbf{Step 1: Uniqueness of the exception terminus.} By BP (\S~\ref{sec:bypass}), $\text{terminus}(\text{ESCALATING}(n)) = \{h(n)\}$. This is not a routing choice — it is an architectural constraint enforced by the bypass mechanism. The escalation payload reaches exactly one entity: the human principal designated at node provisioning. No machine actor may appear in the termination set.

\medskip
\noindent\textbf{Step 2: Binding directives are human-issued only.} The state definitions (\S~\ref{sec:state-machine}) establish that the terminal \texttt{RESOLVED} state is entered only upon issuance of a directive $d \in \{\texttt{ACK}, \texttt{RESOLVE}, \texttt{REJECT}\}$ by $h(n)$. By (C3), no machine actor may issue any of these directives. The \fact{} layer pre-commit hook for the \texttt{RESOLVED} transition verifies the directive's provenance: it checks the cryptographic attestation of the human principal's credential before committing the ledger entry. An unattributed directive cannot enter the ledger.

\medskip
\noindent\textbf{Step 3: Attribution is immutable and non-repudiable.} By (C1), $h(n)$'s credential is established at provisioning and is non-repudiable: the human principal cannot later deny having issued the directive. By (C2), the \fact{} layer records this credential with every directive. By (C4), the Forensic Ledger entry is append-only and tamper-evident — the hash chain ensures that any retrospective alteration is computationally detectable.

\medskip
\noindent\textbf{Step 4: The accountability chain is linear and complete.} For any node $n$, the full accountability chain for a protocol run is:

\[
n \;\xrightarrow{\text{escalation}}\, h(n) \;\xrightarrow{\text{directive}}\, \text{Fact Layer ledger} \;\xrightarrow{\text{hash chain}}\, \text{immutable record}.
\]

Every hop in this chain is deterministic and auditable. The chain contains exactly one human principal at the terminus and exactly one machine actor at the origin. There is no branch point, no alternative terminus, and no intermediate machine-actor absorber.

\medskip
\noindent\textbf{Conclusion.} Combining Steps 1--4, every \texttt{RESOLVED} state in every protocol run is attributable to exactly one named human principal, and no machine actor can issue, simulate, or substitute this attribution. $\square$

% --------------------------------------------------------------
\subsection{Corollary: Unconditional Upstream Accountability}
% --------------------------------------------------------------

The three theorems together establish the upstream accountability guarantee as a formal invariant of the \bxthree{} Framework:

\begin{bluebox}
\textbf{Corollary (Upstream Accountability).} Under the union of assumptions (A1)--(A5) and (L1)--(L3) and (C1)--(C4), for every \bxthree{} node $n$ and every exception condition $x$ that satisfies $\text{TRIGGER}(n, x)$:

\[
\text{ESCALATING}(n, x) \;\implies\; \bigl(\text{terminus} = \{h(n)\}\bigr) \;\land\; \bigl(\text{reachability}(h(n), t \leq T_{\max})\bigr) \;\land\; \bigl(\text{directive attributed to } h(n)\bigr).
\tag{UPSTREAM-GUARANTEE}
\]

In plain terms: every unresolvable exception in any \bxthree{} node \emph{must} propagate to a named human principal, \emph{must} do so within a bounded time, and \emph{must} be recorded as a binding human decision. The guarantee is unconditional — it does not depend on machine-actor cooperation, on the availability of particular resolution pathways, or on the absence of faults in the agent population. It is a structural property of the protocol's architecture.
\end{bluebox}

This corollary is the formal expression of the informal commitment stated in \S~\ref{sec:bailout}: that \bxthree{} systems \emph{fail upward into human consciousness, never downward into autonomous chaos}.
